\documentclass{ximera}
\title{Circles, CW2}


\newcommand{\pskip}{\vskip 0.1 in}

\begin{document}
\begin{abstract}
More questions about circles.
\end{abstract}
\maketitle


\pskip

\section{Equations of Circles}

\begin{question} \label{QDEWW9832}

Let ${\cal C}$ be the circle centered at the point $A$ with coordinates $(-2,1)$ that passes through the point $Q$ with coordinates $(5,-2)$. The coordinates are measured in centimeters.

\begin{enumerate}
\item Sketch this circle. Start by drawing the coordinate axes, labelled with the variable names $x$, $y$, and with units. Then plot the points $A$ and $Q$. Then draw the circle.

\item Find the exact radius of the circle Do not use a calcluator.

\item Find an equation of the circle. Explain your reasoning. 

Start by saying the point $P$ with coordinates $(x,y)$ (measured in cm) lies on the circle if and only if the distance from $P$ to the center $A$ is equal to $\answer{\sqrt{58}}$ cm.

So algebraically, the point $P$ with coordinates $(x,y)$ (measured in cm) lies on the circle with center $A(-2,1)$ cm and radius $\sqrt{58}$ cm if and only if
\[
     \sqrt{\answer{(x+2)^2 + (y-1)^2}}  = \sqrt{\answer{58}}.
\] 

%\item Use alegebra Find the exact coordinates of the points where the line though the points $(-3,1)$ and $(4,5)$ intersects the circle. Start by finding an equation of the line.
%Do not use a calculator.


\end{enumerate}
\end{question}

\begin{question}  \label{QKFr3Dfer33}

\begin{enumerate}

This question is about finding an equation of the smallest circle passing through the points $A(-9,12)$ and $B(7,-2)$.

\item Drag the center $C$ of the circle below to make the circle pass through the points $A$ and $B$.

\begin{onlineOnly}
    \begin{center}
\desmos{d3co2xy08s}{900}{600}
\end{center}
\end{onlineOnly}

\href{https://www.desmos.com/calculator/d3co2xy08s}{141: Smallest Crircle}

\item Change the radius of the circle by dragging the slider in Line 2. Then drag the circle's center so that $A$ and $B$ lie on the circle.

\item Continue to experiment as in part (b). Try to find the smallest circle through $A$ and $B$.

\item Find an equation of the smallest circle through the points $A(-9,12)$ and $B(7,-2)$. Explain your reasoning. 

\end{enumerate}

\end{question}


\section{Closest and Farthest Points on a Circle}


\begin{question}  \label{QKdf49edr4333}
\begin{enumerate}

\item Drag point $R$ in the worksheet below to approximate the coordinates of the point on the circle of radius $5$ centered at the origin that is closest to the point $B$ with coordinates $(-2,3)$. Then mark that position by dragging point $P$ there.

\begin{onlineOnly}
    \begin{center}
\desmos{6yotul255q}{900}{600}
\end{center}
\end{onlineOnly}

\href{https://www.desmos.com/calculator/6yotul255q}{141: Circles Tangent to a Given Crircle}

\item Drag point $R$ to approximate the coordinates of the point on the circle of radius $5$ centered at the origin that is farthest from $B$. Then mark that position by dragging point $Q$ there.

\item Use algebra, not geometry or vectors, to find the exact coordinates of the points in parts (a) and (b). Do \emph{not} use a calculator. Explain your reasoning.

\item Enter the exact coordinates of these points in Lines 19 and 20 in the above worksheet.

\end{enumerate}
\end{question}



\begin{question}  \label{QKdf49e46r4333}
\begin{enumerate}

\item \item Drag point $R$ in the worksheet below to approximate the coordinates of the point on the circle of radius $5$ centered at the origin that is closest to the point $B$ with coordinates $(-2,3)$. Then mark that position by dragging point $P$ there.

\begin{onlineOnly}
    \begin{center}
\desmos{pvomdwhlcd}{900}{600}
\end{center}
\end{onlineOnly}

\href{https://www.desmos.com/calculator/pvomdwhlcd}{141: Circles Tangent to a Given Crircle 8}

\item Drag point $R$ to approximate the coordinates of the point on the circle of radius $5$ centered at the origin that is farthest from $B$. Then mark that position by dragging point $Q$ there.

\item Use algebra, not geometry or vectors, to find the exact coordinates of the points in parts (a) and (b). Do \emph{not} use a calculator. Explain your reasoning.

\item Enter the exact coordinates of these points in Lines 21 and 22 in the above worksheet.
\end{enumerate}
\end{question}



\end{document}